% Part: many-valued-logic
% Chapter: infinite-valued-logics
% Section: introduction

\documentclass[../../../include/open-logic-section]{subfiles}

\begin{document}

\olfileid[pt]{mvl}{inf}{int}

\olsection{Introdução}

O número de valores de verdade de uma matriz não precisa ser finito.
Uma escolha óbvia para um conjunto com infinitos valores de verdade é
o conjunto dos números racionais entre $0$ e~$1$,
$V_\infty = [0,1] \cap \Rat$, isto é,
\begin{align*}
    V_\infty & = \Setabs{\frac{n}{m}}{n,m \in \Nat \text{ e } n\le m}.
\intertext{Ao considerar esse conjunto infinito de valores de verdade,
é frequentemente útil considerar também os subconjuntos}
V_m & = \Setabs{\frac{n}{m-1}}{n \in \Nat \text{ e } n\le m}
\intertext{Por exemplo, $V_5$ é o conjunto com $5$ valores de verdade
igualmente espaçados,}
V_5 & = \{0, \frac{1}{4}, \frac{1}{2}, \frac{3}{4}, 1\}.
\end{align*}
Nas lógicas baseadas nesses conjuntos de valores de verdade, em geral
somente $1$ é designado, isto é, $V^+ = \{1\}$. Em outras palavras,
fazemos $1$ desempenhar o papel da verdade (absoluta) e $0$ o da
falsidade absoluta, mas as !!{formula}s podem assumir qualquer valor
intermediário em~$V$.

Também se pode considerar o conjunto $V_{[0,1]} = [0,1]$ de todos os
números \emph{reais} entre $0$ e~$1$, ou ainda outros subconjuntos
infinitos de $[0,1]$. As lógicas com esse conjunto de valores de
verdade são frequentemente chamadas \emph{difusas}.


\end{document}
